\section{Introduction}

Global investors rebalance toward assets that preserve value when financial markets turn turbulent. \textcite{baur_lucey_2010} define a hedge as a security that is uncorrelated with stocks or bonds on average, and a safe haven as a security that is uncorrelated with stocks and bonds during a market crash. \textcite{hossfeld_macdonald_2014} transpose the same distinction onto currencies, classifying a safe haven currency as one whose effective returns are negatively related to global stock market returns in periods of high financial stress. The Japanese yen and Japanese government bonds (JGBs) are the archetypal case. Foreign investors bought Japanese assets during the 2001 dot-com collapse and the 2007 to 2008 global financial crisis, and the yen has appreciated against the dollar in nearly every major risk-off episode since the 1990s \parencite{beirne_sugandi_2023}. The stakes of this behavior are real. Safe-haven inflows appreciate the currency, erode export competitiveness, and can transmit a global shock into the domestic economy through the exchange rate.

\textcite{beirne_sugandi_2023} provide the most complete accounting of these spillovers for Japan. The replicated article appears in the Pacific-Basin Finance Journal, which carries an A classification on the Australian Business Deans Council Journal Quality List. Their study estimates country-specific structural vector autoregressions (VARs) for Japan, Switzerland, and the United States on working-days data and identifies four results for Japan. The yen's real effective exchange rate (REER) appreciates sharply and persistently after a risk-off shock. Portfolio flows to Japan show no significant response, which the authors attribute to rapid adjustment of financial markets toward new equilibrium exchange rates. Negative real spillovers appear only for Japan, with the appreciation amplifying the decline in GDP growth. Domestic economic policy uncertainty does not respond significantly. The surrounding literature is consistent with these patterns. \textcite{habib_stracca_2012} identify net foreign asset positions as the most robust predictor of safe haven status. \textcite{ranaldo_soderlind_2010} and \textcite{fatum_yamamoto_2015} document the yen's crisis-time appreciation, and \textcite{botman_carvalho_lam_2013} show that yen appreciation during risk-off episodes coincides with offshore derivative trading rather than measurable cross-border flows.

Three developments since the paper's sample ended in March 2021 limit what the published evidence can say today. The first is the regime shift in the yen itself. Between 2021 and 2026 the currency depreciated from roughly 105 to 160 against the dollar despite foreign exchange interventions in 2022 and 2024 that arrested but did not reverse the slide, while the Bank of Japan exited yield curve control and the Federal Reserve raised rates aggressively, an environment in which the traditional safe-haven response weakened. The second is data revision. Japan's real GDP underwent a base-year change and subsequent benchmark revisions that raised current-vintage output by 3.5 to 6.3 percent over 2019 to 2021, and the BIS retroactively rebased its REER from a 2010 base to a 2020 base, so a researcher who downloads these series today does not obtain the data the original authors observed. The third is the set of unreported estimation settings. The paper does not state which quadratic interpolation variant was used, what lag maximum was searched, or how the confidence intervals were computed, and any independent replication therefore carries uncertainty about whether differences originate in data, code, or configuration.

This study replicates \textcite{beirne_sugandi_2023} for Japan and extends the sample through June 2026. The replication rebuilds the paper's recursively restricted VAR with a Cholesky identification scheme, adds the unrestricted specification as a robustness check, and estimates both at the daily and monthly frequencies. The data follow a dual-vintage policy. The paper period from 14 January 1999 to 31 March 2021 uses the 2021-era vintages of real GDP and the REER, and the extension from April 2021 to 30 June 2026 uses current-vintage data. The daily financial market series were downloaded from LSEG Workspace, and the macroeconomic series were collected via automated Python API scripts from the Federal Reserve Bank of St.\ Louis ALFRED archive, the BIS, the Bank of Japan, and the World Uncertainty Index database.

Three research questions guide the analysis. First, \textit{does the paper's qualitative structure for Japan reproduce when the model is re-estimated independently on the original sample}? Second, \textit{do the estimated relationships survive into the post-2021 regime}? Third, \textit{which results are robust to data revisions, and which depend on the exact data vintage}?

Three objectives follow. The first is to \textit{rebuild the restricted and unrestricted SVARs for Japan over the paper period with vintage data and compare the impulse responses against the published figures}. The second is to \textit{re-estimate both specifications over the extended sample and document the differences}. The third is to \textit{disclose every specification choice, including the code defects found and corrected during the replication, so that the boundary between data effects and implementation effects is explicit}.

The study contributes in three ways. It validates the paper's core results for Japan, since the daily restricted VAR reproduces the negative real GDP response, the REER appreciation, the bond spread widening, and the null capital flow responses. It extends the evidence into a regime in which the yen's safe-haven status is disputed, where the real GDP response weakens toward zero. And it demonstrates a data-vintage problem that is material to replication, since the restricted model's GDP response flips sign when current-vintage data replace the 2021-era series.

The remainder of the paper is organized as follows. Section 2 reviews the literature on safe haven currencies and develops the theoretical framework and hypotheses. Section 3 describes the data and the econometric methodology. Section 4 presents the descriptive statistics, stylized facts, and impulse response results. Section 5 discusses the findings in relation to theory and the original paper, and Section 6 concludes. The literature review that follows assembles the currency evidence and derives the five hypotheses that the methodology in Section 3 is designed to test.
